%%
% BIThesis 研究生学位论文模板 The BIThesis Template for Graduate Thesis
% This file has no copyright assigned and is placed in the Public Domain.
%%

\begin{abstract}
  面向垂直领域知识图谱构建的知识抽取任务，既要解决领域标注数据稀缺、自动标注噪声较高的问题，又要兼顾抽取结果的结构可控性与本地部署可行性。针对这一需求，本文以军事新闻领域为实例，围绕“数据构建—可控抽取—图谱应用”三条链路开展研究，提出了一套面向领域知识图谱构建的高质量可控知识抽取方法。首先，面向高质量训练数据获取困难的问题，构建了基于 SO-CoT 自优化与多模型事实评估的银标数据生成流程，通过少量种子样本、自反馈提示优化、事实评估与回流重抽，提升自动标注结果的可解析性、事实性与合规性，最终形成包含 10647 份军事新闻报道和 40565 个实体关系三元组的领域银标数据集。其次，面向轻量化本地部署场景，提出参数高效的本体约束知识抽取方法，以小参数开源模型为基座，结合 QLoRA 完成领域适配，并在推理阶段引入本体约束与多候选结果评分机制，提高抽取结果的结构合法性与稳定性。在 1000 条测试样本上的实验表明，完整方法的 F1 值达到 74.95\%，优于 QLoRA-only 及若干主流实体关系抽取基线，且新增开销主要集中在推理后处理阶段，具备较好的工程可行性。最后，基于上述方法设计并实现了军事新闻领域知识图谱构建与问答原型系统，完成了从文本导入、知识抽取、结果修订、图谱入库到基础查询与简单问答的全流程验证。系统在 1000 条输入文本上构建出 4546 个节点、6017 条关系的知识图谱，验证了本文方法在稳定入库与下游应用中的有效性。研究结果表明，本文方法能够较好地兼顾训练数据质量、抽取结果可控性与本地部署需求，为面向新闻领域知识图谱构建的知识抽取研究提供了一条较完整、可落地的技术路径。
\end{abstract}

% 如需手动控制换行连字符位置，可写 aa\-bb，详见
% https://bithesis.bitnp.net/faq/hyphen.html

\begin{abstractEn}
  Knowledge extraction for vertical-domain knowledge graph construction must address not only the scarcity of domain-specific annotated data and the noise in automatic annotation, but also the controllability of extracted structures and the feasibility of local deployment. To this end, this thesis takes military news as a case study and investigates the problem along three linked stages: data construction, controllable extraction, and knowledge graph application. A high-quality and controllable knowledge extraction method for domain knowledge graph construction is proposed. First, to address the difficulty of obtaining high-quality training data, a silver-standard data construction pipeline based on SO-CoT self-optimization and multi-model factual evaluation is designed. Through a small set of seed samples, self-feedback prompt optimization, factual assessment, and iterative re-extraction, the proposed method improves the parsability, factual consistency, and schema compliance of automatically generated annotations, and finally builds a domain silver-standard dataset containing 10,647 military news reports and 40,565 entity-relation triples. Second, for lightweight local deployment, a parameter-efficient ontology-constrained knowledge extraction method is proposed. A small open-source model is adapted to the target domain with QLoRA, and ontology constraints together with a multi-candidate scoring mechanism are introduced at inference time to improve the structural validity and stability of extracted results. Experiments on 1,000 test samples show that the complete method achieves an F1 score of 74.95, outperforming the QLoRA-only baseline without constraints, while the additional cost is mainly introduced in the post-inference stage, indicating good engineering feasibility. Finally, based on the proposed methods, a military news domain knowledge graph construction and question-answering prototype system is designed and implemented, completing the full workflow from text import, knowledge extraction, and result revision to graph storage, basic query, and simple question answering. On 1,000 input texts, the system constructs a knowledge graph containing 4,546 nodes and 6,017 relations, which verifies the effectiveness of the proposed method in stable graph insertion and downstream application. The results show that the proposed approach can effectively balance training data quality, extraction controllability, and local deployment requirements, and provides a relatively complete and practical technical path for knowledge extraction in news-domain knowledge graph construction.
\end{abstractEn}
